\section*{Chapter \ref{ch:b1:rifts-ridges-margins} --- Rifts, Ridges and Passive Margins}

\begin{solution}{exo:b1:rifts-ridges-margins:1}
(a) The dyke would cover $\qty{60000}{m}\times\qty{8}{m} = \qty{4.8e5}{m^2}$ in
plan, so its mean height would be $\qty{2.5e9}{m^3}/\qty{4.8e5}{m^2} =
\qty{5200}{m}$, about \qty{5}{km}: the dyke cut through much of the \omterm{def:b1:layered-earth:layers}{crust}.
(b) $\qty{8}{m}/\qty{0.018}{m.yr^{-1}} = 444$ years, some four centuries.
\end{solution}

\begin{solution}{exo:b1:rifts-ridges-margins:2}
(a) The \omterm{def:b1:cross-sections:throw}{heave} is the \omterm{def:b1:cross-sections:throw}{throw} divided by the tangent of the \omterm{def:b1:geological-maps:attitude}{dip}:
$3.0/\tan\ang{60} = \qty{1.73}{km}$ and $0.8/\tan\ang{70} = \qty{0.29}{km}$;
the extension of the section is \qty{2.02}{km}. (b) Before the faulting the
section was $25 - 2.02 = \qty{22.98}{km}$ wide, about \qty{23}{km}: it has been
lengthened by about \qty{9}{\percent}.
\end{solution}

\begin{solution}{exo:b1:rifts-ridges-margins:3}
Ultraslow (below \qty{20}{mm.yr^{-1}}): the Gakkel Ridge (12), the central Red
Sea (15) and the Reykjanes Ridge (19), the last close to the limit. Slow (20 to
\qty{55}{mm.yr^{-1}}): the western Carlsberg Ridge (23) and the South Atlantic
(33). Fast (above \qty{75}{mm.yr^{-1}}): the East Pacific Rise (143). The
ultraslow and slow ridges should have an axial valley between normal faults; the
East Pacific Rise an axial high.
\end{solution}

\begin{solution}{exo:b1:rifts-ridges-margins:4}
(a) $\qty{3.0}{km^2}\times\qty{7.1}{km} = \qty{21.3}{km^3}$, about
\qty{21}{km^3} a year. (b) $\qty{3.028e8}{km^2}/\qty{3.0}{km^2.yr^{-1}} =
\qty{1.01e8}{yr}$, \qty{101}{Myr}. The oldest floor, about \qty{190}{Myr}, is
older because the floor is not destroyed in order of age: subduction consumes
floor of every age, some old floor survives where no trench has reached it,
and the rate of production has varied.
\end{solution}

\begin{solution}{exo:b1:rifts-ridges-margins:5}
(a) $2500 + 350\times2 = \qty{3200}{m}$, $2500 + 350\times6 = \qty{4600}{m}$,
$2500 + 350\times8 = \qty{5300}{m}$. (b) $t = \bigl((d - 2500)/350\bigr)^2$:
$(2100/350)^2 = \qty{36}{Myr}$ and $(2800/350)^2 = \qty{64}{Myr}$. (c) The
floor sinks \qty{1400}{m} between 4 and \qty{36}{Myr}, but only \qty{700}{m}
between 36 and \qty{64}{Myr}, over a similar length of time.
\end{solution}

\begin{solution}{exo:b1:rifts-ridges-margins:6}
The effective expansion coefficient is
$\num{3.2e-5}\times3300/2300 = \qty{4.591e-5}{K^{-1}}$, and
$\sqrt{\kappa\times\qty{1}{Myr}} = \sqrt{\num{e-6}\times\num{3.156e13}} =
\qty{5618}{m}$. Then $C = (2/\sqrt\pi)\times\num{4.591e-5}\times1350\times5618 =
1.1284\times\num{4.591e-5}\times1350\times5618 = \qty{393}{m.Myr^{-1/2}}$. For a
\omterm{def:b1:layered-earth:layers}{mantle} \qty{250}{K} colder, $C$ scales with $T_m - T_0$:
$393\times1100/1350 = \qty{320}{m.Myr^{-1/2}}$. The floor of a colder ocean
sinks more slowly and follows the smaller constant; the eastern Pacific, with
\num{350} against \num{390}, is an example.
\end{solution}

\begin{solution}{exo:b1:rifts-ridges-margins:7}
(a) $\qty{e-6}{cal.cm^{-2}.s^{-1}} = \num{4.184e-6}\times\num{e4} =
\qty{4.184e-2}{W.m^{-2}} = \qty{41.84}{mW.m^{-2}}$, so the constant is
$11.3\times41.84 = \qty{473}{mW.m^{-2}.Myr^{1/2}}$, \qty{3.5}{\percent} below
the 490 of \cref{ch:b1:heat-budget}: the two fits agree. (b) $q = 473$,
$473/\sqrt{10} = 150$ and $473/\sqrt{150} = \qty{38.6}{mW.m^{-2}}$. The
square-root law keeps falling, while the plate model levels off at a constant:
at \qty{150}{Myr} the square-root value is too low, as the depth law goes too
deep.
\end{solution}

\begin{solution}{exo:b1:rifts-ridges-margins:8}
Take the northern segment on $x = 0$ and the southern one on $x = -40$
(km), the transform on $y = 0$ between them. (a) Just north of the transform the
floor lies west of the northern segment and moves west; just south of it, it
lies east of the southern segment and moves east. Seen from either side, the
far side moves to the left: the slip is left-lateral (sinistral). (b) The
ridge seems to have been displaced right-laterally, its southern part to the
west: the opposite sense, as \cref{prop:b1:rifts-ridges-margins:transform}
predicts. (c) Earthquakes occur only on the \qty{40}{km} between the two ridge
ends, where the two sides belong to different plates; beyond, along the
\omterm{def:b1:rifts-ridges-margins:transform}{fracture zones}, both sides belong to the same plate and move together.
\end{solution}

\begin{solution}{exo:b1:rifts-ridges-margins:9}
The half rate is \qty{10.6}{mm.yr^{-1}}, \qty{10.6}{km.Myr^{-1}}. (a) The age
step is $350/10.6 = \qty{33.0}{Myr}$. (b) At the end of the transform, floor
\qty{33.0}{Myr} old ($2500 + 350\sqrt{33.0} = \qty{4511}{m}$) faces floor of
age 0 (\qty{2500}{m}): a step of about \qty{2010}{m}. \qty{100}{km} beyond,
the ages are $450/10.6 = 42.5$ and $100/10.6 = \qty{9.4}{Myr}$, the depths
4781 and \qty{3575}{m}: a step of about \qty{1210}{m}. \qty{1000}{km} beyond,
the ages are 127.4 and \qty{94.3}{Myr}, the depths 6450 and \qty{5899}{m}:
\qty{550}{m}. (These ages are beyond \qty{80}{Myr}, where the square-root law
goes too deep; the plate law would give a smaller step.)
\end{solution}

\begin{solution}{exo:b1:rifts-ridges-margins:10}
(a) With equal areas for every age, the mean depth is the mean of $d(t)$ over
$[0, A]$:
\[
  \frac1A\int_0^A\bigl(d_r + C\sqrt t\bigr)\mathrm dt = d_r + \frac CA\cdot\frac23A^{3/2}
  = d_r + \frac23C\sqrt A.
\]
(b) For $A = \qty{180}{Myr}$: $2500 + \frac23\times350\times13.416 =
\qty{5630}{m}$; for \qty{120}{Myr}: $2500 + \frac23\times350\times10.954 =
\qty{5056}{m}$, \qty{574}{m} shallower. (c) Let the sea rise by $r$ and the
loaded floor sink by $s = r\rho_w/(\rho_m - \rho_w)$. The water, unchanged in
volume, must fill the same depth over the same area, so $r + s = \qty{574}{m}$,
$r\,\rho_m/(\rho_m - \rho_w) = 574$, and $r = 574\times2300/3300 =
\qty{400}{m}$: hundreds of metres, from the spreading rate alone.
\end{solution}

\begin{solution}{exo:b1:rifts-ridges-margins:11}
(a) $\sqrt{x - a} = \sqrt x\sqrt{1 - a/x} \approx \sqrt x\,(1 - a/2x)$ for
$a \ll x$, so
$\Delta d \approx C\sqrt{x/u}\cdot a/(2x) = Ca/(2\sqrt{ux})$.
(b) At $x = \qty{450}{km}$: exact \qty{1205}{m}, approximate
$350\times350/(2\sqrt{10.6\times450}) = \qty{887}{m}$; at \qty{1350}{km}:
\qty{550}{m} and \qty{512}{m}; at \qty{3350}{km}: \qty{334}{m} and
\qty{325}{m}. The ratio of the approximate to the exact step depends only on
$r = a/x$: $(r/2)/(1 - \sqrt{1 - r})$, which is 0.9 for $r = 0.36$. The
approximation is good to \qty{10}{\percent} beyond $x = 350/0.36 \approx
\qty{970}{km}$, about three offsets from the ridge.
\end{solution}

\begin{solution}{exo:b1:rifts-ridges-margins:12}
(a) $d_2 - d_1 = 3671.1 - 3606.8 = \qty{64.3}{m}$ at \qty{10}{Myr};
$4072.8 - 4065.3 = \qty{7.5}{m}$ at 20; $4415.4 - 4417.0 = \qty{-1.7}{m}$ at
30; $4707.5 - 4713.6 = \qty{-6.1}{m}$ at 40. (b) $d_2 - d_1$ is positive at 20
and negative at 30, and both laws are continuous: they are equal in between.
Bisection: at 25, $d_2 - d_1 = +0.9$; at 27.5, $-0.7$; at 26.25, $+0.01$; at
26.875, $-0.3$; at 26.56, $-0.2$; at 26.41, $-0.08$; at 26.33, $-0.04$: the
root lies between 26.25 and 26.33, at \qty{26.3}{Myr}, where both give
\qty{4294}{m}. (c) $d_1' = 175/\sqrt{26.3} = \qty{34.1}{m.Myr^{-1}}$ and
$d_2' = (3200/62.8)\,e^{-26.3/62.8} = \qty{33.5}{m.Myr^{-1}}$. The laws have
nearly the same value and the same slope there, so the joined curve has no
visible kink.
\end{solution}

\begin{solution}{pb:b1:rifts-ridges-margins:1}
\textbf{1.} \omterm{def:b1:cross-sections:throw}{Heave} $= $ \omterm{def:b1:cross-sections:throw}{throw}$/\tan$(\omterm{def:b1:geological-maps:attitude}{dip}): F1 \qty{346}{m}, F2 \qty{770}{m},
F3 \qty{2014}{m}, F4 \qty{163}{m}, F5 \qty{462}{m}, F6 \qty{1172}{m}. F3, with
the largest \omterm{def:b1:cross-sections:throw}{throw} and \omterm{def:b1:cross-sections:throw}{heave}, is the border fault.

\textbf{2.} The extension is the sum, \qty{4927}{m}, about \qty{4.9}{km}. The
rift was $60 - 4.9 = \qty{55.1}{km}$ wide before the faulting.

\textbf{3.} The normal to a rift trending \ang{020} points to \ang{110}; the
motion, towards \ang{105}, makes \ang{5} with it. The perpendicular component
is $6.0\cos\ang{5} = \qty{5.98}{mm.yr^{-1}}$.

\textbf{4.} $\qty{4.927}{km}/\qty{5.98}{km.Myr^{-1}} = \qty{0.82}{Myr}$.

\textbf{5.} In \qty{3.2}{Myr} the plates moved apart by
$5.98\times3.2 = \qty{19.1}{km}$ across the rift. The faults carry
$4.93/19.1 = \qty{26}{\percent}$ of it.

\textbf{6.} The other three quarters, \qty{14.2}{km}, are not in the faults of
the section. They may be taken up by dykes intruded along the axis, which open
the rift without faulting it, as in Afar, or by faults outside the section.
With dykes \qty{8}{m} wide, $\qty{14200}{m}/\qty{8}{m} \approx 1780$ dykes are
needed, one every $3.2\times10^6/1780 \approx 1800$ years on average.

\textbf{7.} From C2An: $(21.0 + 20.2)/2.581 = 41.2/2.581 =
\qty{15.96}{mm.yr^{-1}}$. From C3n: $(41.0 + 40.2)/5.235 = 81.2/5.235 =
\qty{15.51}{mm.yr^{-1}}$.

\textbf{8.} $(41.2 - 2.2)/2.581 = \qty{15.11}{mm.yr^{-1}}$ and
$(81.2 - 2.2)/5.235 = \qty{15.09}{mm.yr^{-1}}$: corrected, the two agree to
\qty{0.02}{mm.yr^{-1}}. The correction matters more for the younger chron, whose
distance is shorter.

\textbf{9.} The outward displacement shifts both edges of a flank by the same
amount, so it cancels in their difference. Over both flanks,
$(81.2 - 41.2)/(5.235 - 2.581) = 40.0/2.654 = \qty{15.07}{mm.yr^{-1}}$.

\textbf{10.} North-east: $(41.0 - 21.0)/2.654 = \qty{7.54}{mm.yr^{-1}}$;
south-west: $(40.2 - 20.2)/2.654 = \qty{7.54}{mm.yr^{-1}}$. The spreading is
symmetric; the \qty{0.8}{km} difference between the flanks, the same for both
chrons, is a constant shift, such as the axis being placed \qty{0.4}{km} off
centre.

\textbf{11.} $\cos\Delta = \sin\ang{30.9}\sin\ang{18.0} +
\cos\ang{30.9}\cos\ang{18.0}\cos\ang{16.4} = 0.15870 + 0.78287 = 0.94157$, so
$\Delta = \ang{19.69}$. With $\omega R = 0.403\times(\pi/180)\times6371 =
\qty{44.81}{mm.yr^{-1}}$, the rate is $44.81\sin\ang{19.69} =
\qty{15.10}{mm.yr^{-1}}$, the corrected rates of questions 8 and 9.

\textbf{12.} $\cos\Delta = \sin\ang{30.9}\sin\ang{26.0} +
\cos\ang{30.9}\cos\ang{26.0}\cos\ang{11.4} = 0.22512 + 0.75600 = 0.98112$, so
$\Delta = \ang{11.15}$ and the rate is $44.81\sin\ang{11.15} =
\qty{8.7}{mm.yr^{-1}}$.

\textbf{13.} The predictions, 15.1 and \qty{8.7}{mm.yr^{-1}}, match the
measured 15 and \qty{9}{mm.yr^{-1}}. The rate is $\omega R\sin\Delta$, and the
pole lies north-west of the Red Sea, in north-eastern Libya: the farther south, the farther
from the pole, the faster the opening.

\textbf{14.} $15.09\times5.235 = \qty{79}{km}$ of oceanic floor.

\textbf{15.} $5000/120.6 = \qty{41}{km.Myr^{-1}}$, \qty{41}{mm.yr^{-1}} on
average.

\textbf{16.} $5000/15.1 = \qty{331}{Myr}$ at the central rate;
$5000/8.7 = \qty{575}{Myr}$ at the northern rate.

\textbf{17.} Opening about a nearby pole, the Red Sea widens like a pair of
scissors: its width grows in proportion to $\sin\Delta$, and it stays a wedge
narrowing towards the pole. At \qty{8.7}{mm.yr^{-1}}, the northern part needs
$79/8.7 = \qty{9.1}{Myr}$ to make the \qty{79}{km} of oceanic floor the
central part has now.

\textbf{18.} At \qty{15}{mm.yr^{-1}} the Red Sea axis is an ultraslow ridge; an
axial valley between normal faults is expected.

\textbf{19.} $2500 + 350\times1 = \qty{2850}{m}$;
$2500 + 350\sqrt{10} = \qty{3607}{m}$; $2500 + 350\sqrt{40} = \qty{4714}{m}$.

\textbf{20.} $t = (1500/350)^2 = \qty{18.4}{Myr}$.

\textbf{21.} In the first \qty{10}{Myr} it sinks $350\sqrt{10} = \qty{1107}{m}$;
between 70 and \qty{80}{Myr}, $350(\sqrt{80} - \sqrt{70}) = 350\times(8.944 -
8.367) = \qty{202}{m}$, five times less over the same time.

\textbf{22.} $\mathrm dd/\mathrm dt = C/(2\sqrt t) = 175/\sqrt{80} =
\qty{19.6}{m.Myr^{-1}}$, about \qty{0.02}{mm} a year.

\textbf{23.} $490/\sqrt{80} = \qty{54.8}{mW.m^{-2}}$.

\textbf{24.} Square-root law: $2500 + 350\sqrt{80} = 2500 + 350\times8.944 =
\qty{5630}{m}$. Plate law: $6400 - 3200\,e^{-80/62.8} = 6400 - 3200\times0.2798
= \qty{5505}{m}$. They differ by \qty{126}{m}. At \qty{80}{Myr} both laws still
apply: the square-root law up to about \qty{80}{Myr}, the plate law from about
\qty{20}{Myr}; beyond, the square-root law goes too deep.

\textbf{25.} The square-root law gives \qty{5630}{m}: to the nearest tenth of a
kilometre, the floor of the new ocean lies at a depth of
\textbf{\qty[text-series-to-math=true]{5.6}{km}} when it is 80 million years
old (the plate law gives \qty{5.5}{km}).
\end{solution}
